\section{The connected logarithm}
\label{subsec:rgnl-cluster-logarithm}

Two polymers are incompatible when their union is connected
through faces. In particular a polymer is incompatible with
itself. For an ordered tuple $(X_1,\ldots,X_n)$ define
$u^T(X_1)=1$ and
\[
 u^T(X_1,\ldots,X_n)=\sum_{G\text{ connected}}
                                      (-1)^{|G|},\quad n\ge2,
\]
where $G$ ranges over spanning subgraphs of the incompatibility
graph of the tuple. The coefficient is zero when that graph
is disconnected.

We will use the formal connected-graph identity
\begin{equation}\label{eq:rgnl-connected-graph-identity}
 \sum_{\{X_i\}\text{ compatible}}\prod_i z_{X_i}
 =\exp\left\{\sum_{n\ge1}\frac1{n!}
       \sum_{(X_1,\ldots,X_n)}u^T(X_1,\ldots,X_n)
                                           \prod_i z_{X_i}\right\}
\end{equation}
and the tree bound
\begin{equation}\label{eq:rgnl-tree-domination}
 |u^T(X_1,\ldots,X_n)|
       \le\#\{\text{spanning trees of the incompatibility graph}\}.
\end{equation}
Here the identity is one of formal series in commuting
variables, before substitution into a normed algebra.
For its proof, express the left side as an ordered sum with
factor $1/n!$ and multiply by
$\prod_{i<j}(1-\mathbf1_{\{X_i,X_j\}\text{ incompatible}})$.
Expanding this product sums over all graphs. A graph is
uniquely the union of its connected components. Grouping by
their vertex sets gives the exponential on the right:
an ordered list of component sizes $n_1,\ldots,n_r$ has
$n!/(n_1!\cdots n_r!)$ vertex assignments, and division
by $r!$ removes the order of its components.
Only finitely many terms contribute to a fixed degree.

For \eqref{eq:rgnl-tree-domination}, let $\mathbf G$ be
the incompatibility graph and suppose it is connected.
For a connected spanning subgraph $G$, root at $1$ and,
for $i\ne1$, choose as its parent the least labelled
neighbor at distance $d_G(i)-1$ from the root.
These edges form a tree $T(G)$ with the same root distances
as $G$. Given a spanning tree $\tau$, let $R(\tau)$ contain
its edges, all edges of $\mathbf G$ within one distance
level, and all edges from $i$ to the preceding level whose
other endpoint exceeds the parent of $i$ in $\tau$.
Then
\[
 T(G)=\tau\quad\Longleftrightarrow\quad
                                      \tau\subseteq G\subseteq R(\tau).
\]
In one direction the distances and least parents force
every edge of $G$ to belong to $R(\tau)$.
Conversely, along any edge of $R(\tau)$ the tree distance
changes by at most one. Thus a graph between $\tau$ and
$R(\tau)$ has the tree's distances, and the least neighbor
in the preceding level is its tree parent.
The connected subgraphs are therefore partitioned into
these intervals, and their signed sum is
\[
 \sum_\tau(-1)^{n-1}
           \sum_{A\subseteq R(\tau)\setminus\tau}(-1)^{|A|}
 =(-1)^{n-1}\#\{\tau:R(\tau)=\tau\}.
\]
This proves the tree bound, including the cases
of a disconnected graph and $n=1$.

For this section and the symmetry proof, write $\widehat\mu=\widehat\mu_k$ for the normalized product box law in \eqref{eq:rg-step-box-normalization}.

\begin{proposition}\label{prop:rgnl-cluster-logarithm}
Let $V(X)$ be even elements in the old and fluctuation
generators, continuous on $\mathrm{Box}_k$, with
\[
 \|V(X)\|_{(w;\mathsf w_W)}
                       \le H_{\mathrm{cl}}\exp(-a_2d_M(X)).
\]
Assume that $V(X)$ has fluctuation indices only in $X$ and
depends on $\zeta$ only through coordinates supported in $X$.
For any union $Y$ of $M$-cubes put
\begin{align}
 K(Y)&=\sum_{\substack{\{X_i\}\text{ distinct}\\\bigcup_iX_i=Y}}
                    \prod_i(\exp V(X_i)-1),\qquad K(\varnothing)=1,
                                      \label{eq:rgnl-Mayer-activity}\\
 H(Y)&=\int\dd\mu_I(W)\,\dd\widehat\mu(\zeta)\,K(Y).
                                      \label{eq:rgnl-integrated-activity}
\end{align}
For a connected $Y$ define
\begin{equation}\label{eq:rgnl-log-activity}
 E^\#(Y)=\sum_{n\ge1}\frac1{n!}
   \sum_{\bigcup_iY_i=Y}u^T(Y_1,\ldots,Y_n)\prod_iH(Y_i).
\end{equation}
The last series converges absolutely in the $w$ norm and
\begin{align}
 \int\dd\mu_I(W)\,\dd\widehat\mu(\zeta)
                         \exp\left\{\sum_XV(X)\right\}
       &=\exp\left\{\sum_YE^\#(Y)\right\},
                                      \label{eq:rgnl-log-identity}\\
 \|E^\#(Y)\|_w
       &\le c_S'H_{\mathrm{cl}}\exp(-a_3d_M(Y)),
                                      \label{eq:rgnl-log-bound}
\end{align}
Here $a_3=\kappa-5\kappa_0-4$ and
$c_S'=14c_X^2K_0^2\exp(3\kappa)$.
The sum of the norms of the terms with $n>N$ in
\eqref{eq:rgnl-log-activity} is at most $2^{-N}$ times the
right side of \eqref{eq:rgnl-log-bound}.
\end{proposition}

\begin{proof}
Set
\[
 \begin{gathered}
 a_4=\kappa-3\kappa_0-3,\quad
 a_5=\kappa-3\kappa_0-2,\quad
 a_6=\kappa-5\kappa_0-3,\\
 H_1=7c_XK_0\exp(3a_5)H_{\mathrm{cl}},\qquad
 z_S=28H_1\exp(\kappa_0)K_0\exp(3a_6).
 \end{gathered}
\]
The fixed small-charge choice
\[
 c_S e^{1/4-5\varepsilon}\le1,\qquad
 c_S=10^5c_X^5K_0^5\exp(6\kappa)
\]
implies
\begin{equation}\label{eq:rgnl-cluster-smallness}
 H_{\mathrm{cl}}\le1,\qquad
 4H_{\mathrm{cl}}K_0c_X\exp(3a_5)\le1,\qquad
                                      z_S\le\tfrac12.
\end{equation}
Indeed the second left side is
$144c_X^3K_0^3\exp(3a_5)e^{1/4-5\varepsilon}$.
For the third,
$2z_S=14112c_X^3K_0^4
\exp(3a_5+3a_6+\kappa_0)e^{1/4-5\varepsilon}$, and
$3a_5+3a_6+\kappa_0\le6\kappa$.

Even elements commute, so expanding their exponential
gives $\exp(\sum_XV(X))=\sum_YK(Y)$.
Moreover $K(Y)$ factors over the connected components
of $Y$: every polymer in its defining set belongs to
exactly one component.
If distinct polymers $X_1,\ldots,X_n$ have connected
union $Y$, their incompatibility graph is connected.
Otherwise two nonempty groups would partition their
union; connectedness of $Y$ would give coincident or
adjacent cubes in different groups and hence an edge
between them.
For each incompatibility-tree edge, choose coincident or
face-adjacent cubes in the two polymers and points
where their nets meet these cubes. The two points
lie in a box of side lengths $M,M,2M$, so they
can be joined by a segment of length at most
$\sqrt6M<3M$. Joining near-minimal nets along
the tree proves
\[
 d_M(Y)\le\sum_i d_M(X_i)+3(n-1).
\]
Using $\|\exp V-1\|\le2\|V\|$ and $a_2=a_5+\kappa_0$
therefore bounds the product by
\[
 \exp(-a_5d_M(Y))
       \big(2H_{\mathrm{cl}}\exp(3a_5)\big)^n
                                  \prod_i \exp(-\kappa_0d_M(X_i)).
\]
Replace distinct sets by ordered tuples with factor $1/n!$,
and use the counting bound for subpolymers of $Y$. The result is
\[
 \|K(Y)\|\le \exp(-a_5d_M(Y))x_Y\exp(x_Y),\qquad
 x_Y=2H_{\mathrm{cl}}K_0\exp(3a_5)|Y|_M.
\]
By \eqref{eq:rgnl-cluster-smallness},
\[
 \begin{gathered}
 x_Y\le(d_M(Y)+1)/2,\\
 x_Y\le4H_{\mathrm{cl}}K_0c_X\exp(3a_5)\exp(d_M(Y)/2).
 \end{gathered}
\]
Since $4\exp(1/2)\le7$ and $a_5-1=a_4$, we obtain
\begin{equation}\label{eq:rgnl-Mayer-bound}
 \|K(Y)\|_{(w;\mathsf w_W)}
                         \le H_1\exp(-a_4d_M(Y))=:\mu(Y).
\end{equation}

Weighted Berezin integration costs no norm here:
each paired moment at level zero is
$\pm48^{|J|/2}$, which is exactly covered by the
weight $(\sqrt{48})^{|J|}$ of its generators.
The corresponding moment is $\pm1$ at positive levels.
The photon measure is a probability measure on its compact
box. Hence $\|H(Y)\|_w\le\mu(Y)$.
These integrals produce even elements, since a nonzero
Berezin moment uses an even number of generators.
For compatible polymers the fluctuation modes and
supported photon coordinates are disjoint. The two
product measures therefore factor, with the external
old generators left as coefficients. It follows that
\begin{equation}\label{eq:rgnl-compatible-partition}
 \int\dd\mu_I(W)\,\dd\widehat\mu(\zeta)e^{\sum_XV(X)}
             =\sum_{\{Y_i\}\text{ compatible}}\prod_iH(Y_i).
\end{equation}

We now bound the infinite connected series. Let
\[
 S_n(Y)=\sum_{\bigcup_iY_i=Y}
                   |u^T(Y_1,\ldots,Y_n)|\prod_i\mu(Y_i).
\]
For $n\ge2$, apply \eqref{eq:rgnl-tree-domination}.
For every contributing tree the union estimate and
$a_4=a_6+2\kappa_0$ give
\[
 \prod_i \exp(-a_4d_M(Y_i))
 \le \exp(-a_6d_M(Y))\exp(3a_6(n-1))
                                  \prod_i \exp(-2\kappa_0d_M(Y_i)).
\]
Drop the prescribed union, retaining $Y_i\subseteq Y$,
and root the tree at $1$. Write $q_j$ for the number
of children at $j$. A polymer incompatible with $Y'$
contains one of at most $7|Y'|_M$ cubes consisting of
its cubes and their face neighbors. Since
$|Z|_M\le c_X(d_M(Z)+1)\le\kappa_0(d_M(Z)+1)$ and
$a^q\le q!\exp(a)$, the counting estimate gives
\begin{align*}
 \sum_{Z\text{ incompatible with }Y'}
      \exp(-2\kappa_0d_M(Z))|Z|_M^{q}
            &\le7q!\exp(\kappa_0)K_0|Y'|_M,\\
 \sum_{Z\subseteq Y}\exp(-2\kappa_0d_M(Z))|Z|_M^q
            &\le q!\exp(\kappa_0)K_0|Y|_M.
\end{align*}
Summing from leaves to root thus costs at most
$(7\exp(\kappa_0)K_0)^n|Y|_M\prod_jq_j!$.
For prescribed children counts, Cayley's degree formula
\cite[Proposition~4]{PostnikovLin2019}
counts
$(n-2)!/((q_1-1)!\prod_{j\ge2}q_j!)$ trees, where
$q_1\ge1$ and $\sum_jq_j=n-1$.
Consequently
\[
 \sum_{\text{trees}}\prod_jq_j!
   =(n-2)!\sum_{(q_j)}q_1
   \le(n-1)!\binom{2n-2}{n-1}
   \le(n-1)!4^{n-1}.
\]
Collecting the factors, also for $n=1$, gives
\begin{equation}\label{eq:rgnl-connected-n-bound}
 \frac{S_n(Y)}{n!}
 \le7H_1\exp(\kappa_0)K_0|Y|_M
                                      \exp(-a_6d_M(Y))z_S^{n-1}.
\end{equation}
Since $z_S\le1/2$, sum the geometric series and use
$|Y|_M\le c_X\exp(d_M(Y))$ and $a_6-1=a_3$.
The bound becomes
\[
 98c_X^2K_0^2\exp(\kappa_0+3a_5)
                                 H_{\mathrm{cl}}\exp(-a_3d_M(Y))
 \le c_S'H_{\mathrm{cl}}\exp(-a_3d_M(Y)),
\]
because $\kappa_0+3a_5=3\kappa-8\kappa_0-6$.
Summing only $n>N$ gives the additional factor $2^{-N}$.
This proves absolute convergence and the claimed estimates.

It remains to justify the logarithm identity after
substitution in \eqref{eq:rgnl-connected-graph-identity}.
The even subalgebra with the coefficient norm is a
commutative Banach algebra, and the polymer set is finite.
The total norm sum of the connected terms is finite by
\eqref{eq:rgnl-connected-n-bound}.
In the exponential of their sum, the norm sum of all
products of connected terms is bounded by the exponential
of that total norm sum. Regrouping is therefore permitted.
Regroup by monomials in the formal variables $z_Y$.
The coefficients are those of
\eqref{eq:rgnl-connected-graph-identity}, namely one
for compatible sets and zero for all other monomials.
Its substitution $z_Y=H(Y)$ gives
\eqref{eq:rgnl-compatible-partition}. Grouping the
connected terms by their union gives
$\sum_YE^\#(Y)$ and proves \eqref{eq:rgnl-log-identity}.
\end{proof}

\subsection{Polynomial approximation at fixed volume}\label{subsec:rgnl-polynomial-approximation}

There are also polynomial approximants to this logarithm
which converge uniformly for all families satisfying the
preceding bound, at fixed side. Replace $\exp V-1$ by
$p_N(V)=\sum_{1\le j\le N}V^j/j!$ in $K,H$, and stop
the connected series at $n=N$; denote the result $E_N^\#$.
It is a polynomial in integrals of products of $V(X_i)$,
$X_i\subseteq Y$, and has the same bound
\eqref{eq:rgnl-log-bound}. To verify uniform convergence,
write $D=|\mathcal P_k|$. Since $\|V(X)\|\le1$,
\[
 \|\exp V(X)-1-p_N(V(X))\|
                                  \le\frac{\exp(1)}{(N+1)!}.
\]
Telescoping each finite product in $K$ gives
\[
 \|H(Y)-H_N(Y)\|\le
                             \frac{4^DD\exp(1)}{(N+1)!}.
\]
Thus this difference is at most $\delta_N\mu(Y)$, where
\[
 \delta_N=\frac{4^DD\exp(1)}
                      {(N+1)!\min_{Y\in\mathcal P_k}\mu(Y)}
                                                   \longrightarrow0.
\]
Replacing the $n$ factors in a connected term costs at
most $n\delta_N\prod_i\mu(Y_i)$. Summing
\eqref{eq:rgnl-connected-n-bound} with this factor, using
$\sum_{n\ge1}n2^{-(n-1)}=4$, and adding the tail yields
\[
 \sup_Y\|E^\#(Y)-E_N^\#(Y)\|
 \le28\delta_NH_1\exp(\kappa_0)K_0(\ell/M)^3
                                  +2^{-N}c_S'H_{\mathrm{cl}}.
\]
This convergence is uniform in the connections and the
allowed data. The displayed error depends on the side;
no volume-independent approximation is claimed.

